\documentclass[../main]{subfiles}
\begin{document}
\ifSubfilesClassLoaded{\setcounter{page}{4}}{}
\section{Вигадуємо «окуляри-велосипед»}
\label{sec:01_izobretaem_ochki}

Що таке комунізм? Звідки він взявся, куди закликає, кого закликає? На всі ці прості запитання, як показала історія, потрібні відповіді не раз і назавжди, а кожного разу, коли вони виникають. Як сказав сценарист Григорій Горін устами «того самого Мюнхгаузена», «любов – це теорема, яку потрібно доводити щодня». І в нас так само, тільки комунізм буде теоремою (навіть теорією!) ще серйознішою та глобальнішою, ніж любов. І потрібно навчитися спочатку правильно записувати в пункті «дано» для цієї теореми все дійсно існуюче, а не уявне.

Ті претензії, що висуваються до самого слова «комунізм», зумовлені історичним відкатом, який відбувся в нашому суспільстві. Зниження до вульгарності і навіть перетворення цього поняття на лайливе слово після розпаду СРСР сталося через те, що перемогли суспільні сили, відкрито ворожі до комунізму.

До розпаду Радянського Союзу вже почалося «розмивання» цього поняття, і тому деякі теоретики говорили про необхідність розмежування всередині радянської економіки та ідеології, проведення чіткої межі, щоб не сприймати капіталістичну інерцію як комуністичну тенденцію. Після цього – через те, що суспільство було втягнуте в реставрацію капіталізму в умовах майже втраченої традиції комуністичної теоретичної думки, коли питання полягало у збереженні хоч якоїсь її частини – проблеми комунізму та руху до нього спеціально не досліджувалися. Уся увага була зосереджена на розквітаючому капіталізмі. І його прихильники, і його противники говорили про нього як про щось реальне, тоді як комунізм і рух до нього сприймалися як щось із минулого, або, у випадку з прихильниками, як дуже віддалене майбутнє, до якого ще треба дожити, тому важливо зосередитися на тому, щоб дожити. Комунізм як теоретична проблема на десятиліття зник з порядку денного.

У масовій культурі, далекій від теоретичного мислення, завжди активно відтворювалося (і відтворюється зараз) свідоміння правлячого класу. Цей клас боїться комунізму і ненавидить його, втілюючи свій страх і ненависть у форму різних соціальних міфів, на яких у постсоветських країнах зростає вже друге покоління. Однак міфологізовані уявлення про комунізм не мають нічого спільного з поняттям комунізму. Тут ми знаходимо першу найважливішу точку «ворожості» (яка насправді такою не є). Громадяни країн, які перейшли до капіталізму не через інтервенцію, а «власними силами», просто вже не знають, що таке справжній комунізм, тобто науковий, в чому суть ідеї комунізму. І, тим не менш, проблема комунізму знову постає перед ними як практична. І ця цілісна теорія, зневажання та ігнорування якої відіграли не останню роль у ковзанні назад до капіталізму в СРСР та інших країнах соціалістичного табору, потребує відродження та пояснення «простими словами» людям, чиї справжні інтереси вона відображає. Йдеться не лише про політично нейтральних робітників-пролетаріїв, але й про тих, хто сам себе вважає комуністами або принаймні лівими, – тобто про тих, хто політично визначився як симпатизуючий комунізму. Останні сприймають комунізм як певну «хорошу» суспільну систему, трактуючи цю «хорошу» систему скоріше на рівні чуттєвих уявлень, часто «від зворотного» – як відсутність мерзостей і жахів сучасного капіталістичного суспільства. Але позитивної відповіді на питання «що таке комунізм?» такі прихильники не можуть дати навіть собі.

Така ситуація спонукає саме зараз створити цю брошуру як свого роду «окуляри-велосипед», які допоможуть знову побачити комунізм і зрозуміти, що це таке. Щодо противників комунізму, то, принаймні, їх потрібно змусити висувати не дурні, а розумні претензії до вже наукового комунізму.

У рік сторіччя Жовтневої революції в публічному просторі було представлено настільки багато різних інтерпретацій ролі та поведінки більшовиків у 1917 році, що думки та оцінки, як і слід було очікувати, поляризувалися: від благородних героїв, які «врятували державу», до бандитів, які ненавиділи цю державу (державність) і знищили її. Однак, головне, що вислизнуло з уваги під час такої поляризації: більшовики, і, звісно, Ленін, були мислителями, і у своїх діях керувалися не лише бажаннями, уподобаннями тощо, а й теорією, яка стала підсумком всього попереднього розвитку людської думки. Вороги соціалізму та комунізму люблять називати Жовтень і подальше соціалістичне будівництво експериментом, проведеним більшовиками над народами… І ці вороги значно ближче до істини, ніж далекі від теорії прихильники.

І саме Велику Жовтневу революцію можна назвати експериментом, якщо, звісно, не розуміти це слово в банальному сенсі, як зовнішні маніпуляції. Більше ніяк її не можна було назвати, враховуючи її безпрецедентність. Це був експеримент, проведений відповідно до всіх наукових традицій, адже не буває експериментів без гіпотез. І така гіпотеза, розроблена Леніним та його прихильниками і заснована на спадщині Маркса та Енгельса, полягала в наступному: Світова революція. Почавшись у Росії, пролетарська революція мала охопити промислово розвинені європейські країни, давши поштовх для перетворення їх на Всесвітню республіку Рад. Не всі наукові гіпотези підтверджуються експериментом у первісній формі (хоча ідея актуальності пролетарських революцій загалом залишається вірною і зараз), однак і цей експериментальний рух уздовж навіть частково непідтвердженої гіпотетичної траєкторії сприяє прогресу науки (соціалізму).

Запропонована метафора лише повертає нас до іншого способу бачення революції, військового комунізму, нової економічної політики, індустріалізації тощо, і спонукає розглядати сучасне суспільство та тенденції його розвитку крізь призму комунізму як найважливішої науки в системі всіх суспільних наук, необхідних для того, щоб можна було діяти розсудливо в сучасному світі. Щоб революційний клас міг подолати підпорядкування постійній експлуатації, подолати в собі пригноблений клас капіталістичного суспільства і стати на власну класову позицію, усвідомити та розсудливо змінити суспільні відносини в напрямку ліквідації будь-якого пригноблення та експлуатації людини людиною, в напрямку комунізму.

Тому «теорії-велосипеди», висміяні Маяковським в іншому контексті, допоможуть нам лише в тому випадку, якщо ми проявимо наполегливість і «винаходитимемо» їх заново (відтворюватимемо): пройдемо марксистським шляхом до комуни та Інтернаціоналу, ленінським теоретичним шляхом до Жовтневої революції, а потім і всім радянським шляхом, не змінюючи «засоби пересування» – революційну теорію.

«Заново» означає зробити своїм надбанням науку комунізму в сучасних, успадкованих нами, історичних умовах.
\end{document}
